\begin{statement}{B.30}{Proved}{Full relative compatibility over rationals and at primes at least five}
All groups in this statement are over $\mathbb Q$, with their
lower-central inverse-limit topology. Let $\mathcal P_r$ be the
group of outer automorphisms of the free prounipotent fundamental
group of an $(r+1)$-punctured sphere which act trivially on
abelianization and preserve each peripheral conjugacy class.
Use the eight-puncture source and both rank-two targets of B.25.
The actual deck action defines an $S_3$-action on $\mathcal P_7$
and on the target belonging to $d$; identify the latter target
with the absolute target through $\iota_U$. There is a surjective
continuous filtered homomorphism
\[
 \mathcal R:\mathcal P_7^{S_3}
       \longrightarrow\mathcal P_2\times\mathcal P_2^{S_3},
 \qquad [a]\longmapsto(c_*[a],d_*[a]).
\]
Its degree-$n$ map is exactly $R_n$ of B.26. In particular every
$q\in\mathcal P_2^{S_3}$ has a full filtered relative lift with
induced pair $(q,q)$. If $q\ne1$, it also has a lift with pair
$(q,1)$, which is incompatible with the diagonal condition.
These conclusions include rational filtered centralizer elements
with leading term $q_7$ and nonzero pentagon defect. They do not
specify the relative action induced by a profinite word.

There is also an integral pro-$\ell$ version for every
$\ell\ge5$. Replace $U(N),U(K),U(F_0)$ by their free
pro-$\ell$ completions and the relative groups by the
quotients of $\widehat F_0,\widehat B$ retaining those
completed open kernels and their quotient $S_3$.
With $\mathcal P_{r,\ell}$ the corresponding special
outer IA group, the actual map
\[
 \mathcal R_\ell:\mathcal P_{7,\ell}^{S_3}
       \longrightarrow
       \mathcal P_{2,\ell}\times\mathcal P_{2,\ell}^{S_3}
\]
is a surjective continuous homomorphism of profinite groups,
surjective on every lower-central filtered piece. In
particular each pro-$\ell$ centralizer element constructed
in B.12 has a compatible relative pro-$\ell$ lift.
\end{statement}
\begin{proof}
We first work with the actual relative completions, rather than
assuming a simultaneous linearization of the two maps. An extension
of a finite group $H$ by a prounipotent group over $\mathbb Q$
splits, and its splittings are conjugate by the radical. Indeed,
on each successive central vector-group layer its obstruction
is in $H^2(H,V)$ and its conjugacy ambiguity is in $H^1(H,V)$.
Both vanish by averaging by $|H|^{-1}$. Lift a splitting and its
correction one layer at a time and take the inverse limit. All
corrections are rational. Choose a splitting
$j:S_3\to\mathscr R_v(F_0)$ and its image
$\Phi j:S_3\to\mathscr R_p(B)$. These give genuine deck
actions on $U(N)$ and $U(K)$, and $U(d)$ is equivariant.
The induced outer actions are the original deck actions and do
not depend on this choice. Their degree-one actions are those
in B.25.

The quotient of $U(N)$ by the closed normal subgroup generated
by
\[
                 j(g)u j(g)^{-1}u^{-1}
       \quad(g\in S_3,\ u\in U(N))
\]
is exactly $U(F_0)$, with quotient map $U(c)$. To prove this
assertion, a homomorphism from that quotient to any unipotent
$\mathbb Q$-group is the same as a homomorphism from
$U(N)\rtimes j(S_3)$ to that group which kills $j(S_3)$.
Every homomorphism from the relative completion to a unipotent
group kills this finite subgroup, since a characteristic-zero
unipotent group has no nontrivial element of finite order.
The universal properties of relative and absolute completion
therefore identify the quotient with the absolute completion
of $F_0$. This proves the claim as a morphism of proalgebraic
groups, not merely on associated gradeds.

The kernel of $U(d)$ is the closed normal subgroup generated
by the five peripheral classes labelled $A_0,A_1,A_2,C_0,C_1$.
This follows from the filling presentation in B.25: imposing
those peripheral elements as relations gives the free group
of the remaining three-punctured sphere, and prounipotent
completion preserves this presented quotient by its universal
property.

An $S_3$-invariant outer class in $\mathcal P_7$ has a
representative which strictly commutes with the chosen genuine
deck action. For precision, choose a representative $a$ and
put $z_g=a^{-1}g(a)$. Since its outer class is invariant,
$z_g$ is inner; moreover
$z_{gh}=z_g\,g(z_h)$. The group of inner automorphisms of
$U(N)$ is prounipotent. Its finite-group $H^1$ is trivial by
successive averaging on its lower-central vector layers,
as above. Thus $z_g=b^{-1}g(b)$ for an inner automorphism
$b$, and $ab^{-1}$ commutes with the deck action. If the
outer class begins in degree $n$, choose the initial
representative which is the identity through that degree;
the same argument in the corresponding filtered inner
subgroup preserves this condition. The free Lie algebra has
zero centre, so an inner automorphism whose first nonzero
term has degree at least $n$ is induced by an element of
the corresponding lower-central subgroup. This justifies
the filtered assertion.

A strict deck-centralizing representative preserves the
kernel of $c$, by its displayed normal generators. It also
preserves the kernel of $d$, since it preserves all eight
peripheral conjugacy classes. Its inverse does the same;
it therefore descends to automorphisms of both targets.
For $c$, a peripheral element maps to the square or cube
of a target peripheral element. Unique divisibility in a
prounipotent group shows that preservation of the conjugacy
class of that power preserves the class of its root.
For $d$, the surviving three peripheral elements map to
the target peripheral elements themselves. Both descendants
are thus peripheral and act trivially on abelianization;
the second commutes with its deck action. Changing the
representative by an inner automorphism changes the two
descendants by their corresponding inner automorphisms.
This defines the homomorphism $\mathcal R$ on outer classes.
It is filtered and continuous because $c$ and $d$ preserve
the lower-central series.

We now identify its actual filtration. The correction proof
of B.11 works over $\mathbb Q$ for every rank: the bracket
map onto the next homogeneous free-Lie layer cancels the
peripheral product error. Corrections converge in the
lower-central topology. The resulting endomorphism induces
the identity on every associated graded layer, hence is an
automorphism by induction on the nilpotent quotients. It
follows that
\[
           \operatorname{gr}_n\mathcal P_r
                         =\mathfrak s_{r,n}(\mathbb Q).
\]
These are equalities of the actual leading-term spaces,
not just tangent-space inclusions. The outer filtrations
are complete: compatible outer representatives can be made
compatible as genuine representatives by successively
lifting and removing their inner discrepancies. Their
inverse limit is a continuous peripheral automorphism.
Here peripheral conjugators also converge: for a correction
beginning in degree $n$, remove a power of the corresponding
peripheral generator from its conjugator. The centralizer
calculation in the free Lie algebra then puts that conjugator
in degree at least $n$. The products of these correcting
conjugators are Cauchy. Thus preservation of peripheral
classes in the limit is established by actual conjugators.

Taking finite-group invariants commutes with passage to
each of these graded quotients. To see surjectivity, lift
an invariant leading term arbitrarily. Its discrepancies
form a cocycle in the next filtered subgroup of
$\mathcal P_r$. Successive averaging on that subgroup's
rational vector layers makes the lift invariant without
changing its leading term. Thus
\[
 \operatorname{gr}_n(\mathcal P_7^{S_3})
       =\mathfrak s_{7,n}^{S_3},\qquad
 \operatorname{gr}_n(\mathcal P_2^{S_3})
       =\mathfrak s_{2,n}^{S_3}.
\]
The leading-term maps of the actual $c,d$ are exactly the
maps of B.25, so $\operatorname{gr}_n\mathcal R=R_n$.
B.26 proves that each is onto.

Finally let $(a,b)$ be any target pair. Starting with the
identity source element, suppose its image agrees with
$(a,b)$ through degree $n-1$. The remaining degree-$n$
discrepancy has a preimage under $R_n$. Choose an invariant
source element with that leading term and multiply by this
correction on the appropriate side. The image now agrees
through degree $n$. Completeness makes these successive
corrections converge to a source element with image
$(a,b)$. This proves the claimed surjectivity. Choosing
$(q,q)$ and $(q,1)$ proves both lift assertions. The same
rational special and invariant lifting gives an actual
$q$ with leading term $q_7$; its pentagon defect retains
the nonzero first coefficient $16$ of B.12.

For the pro-$\ell$ assertion write
\[
 \begin{aligned}
 R_{v,\ell}&=\widehat F_0/
        \ker(\widehat N\longrightarrow N^{(\ell)}),\\
 R_{p,\ell}&=\widehat B/
        \ker(\widehat K\longrightarrow K^{(\ell)}).
 \end{aligned}
\]
The kernels in these formulas are characteristic in the
open normal subgroups, hence normal in the ambient groups.
Finite index identifies $\widehat N$ and $\widehat K$
with their respective open subgroups. Both displayed
quotients have pro-$\ell$ radical and quotient $S_3$;
$\varphi$ descends because its restriction sends the first
kernel into the second. Profinite Schur--Zassenhaus gives
a complement in $R_{v,\ell}$; its image is a complement
in $R_{p,\ell}$. One can prove the needed assertion by
splitting every finite extension by an $\ell$-group and
using conjugacy of its complements and compactness. The
same argument gives trivial nonabelian $H^1(S_3,H)$ for
every pro-$\ell$ group $H$ with continuous $S_3$-action.

Every homomorphism from $\widehat F_0$ to a pro-$\ell$
group factors through $R_{v,\ell}$ by restriction to
$\widehat N$, and kills its order-six complement. Hence
the quotient of $N^{(\ell)}$ by its deck commutators is
exactly $F_0^{(\ell)}$. The filling kernel is the closed
normal subgroup of the same five peripheral classes, by
the free pro-$\ell$ presentation. These give the two
actual quotient maps. Square and cube maps have unique
inverses on a pro-$\ell$ group when $\ell\ge5$: on each
finite quotient their inverses are powers by the inverse
of $2$ or $3$ modulo its exponent. Thus the preceding
proof that both descendants are peripheral applies.

The special outer IA and filtered inner groups are
closed subquotients of the pro-$\ell$ kernels of the
actions on Frattini quotients. Coprime $H^1$ therefore gives
strict deck-centralizing representatives and surjectivity
onto invariant leading terms. B.11 supplies the actual
integral leading terms. The formulas proving B.26 use
only division by $2$ and $3$, including the cabling and
both averages. They give surjectivity over $\mathbb Z_\ell$
in every degree. Successive corrections now converge in
the actual pro-$\ell$ topology, since every finite
$\ell$-group quotient is nilpotent. Equivalently,
compactness supplies compatible peripheral conjugators
and outer representatives. This proves the pro-$\ell$
surjectivity and its filtered version. No claim at
$\ell=2$ or $3$ is made by this relative argument:
its complement and averaging steps require $6$ to be
invertible.

No automorphism of the full free profinite group has been
constructed in this argument. For such an action, the source relative action
is the particular one induced on the open subgroup $N$;
here it was chosen freely to solve two projection equations.
The existence of two such actions with the same first
projection demonstrates why these full rational and relative
pro-$\ell$ results still do not evaluate the prescribed word in B.28.
\end{proof}
